\documentclass[11pt]{article}
\usepackage[a4paper,margin=25mm]{geometry}
\usepackage{amsmath,amssymb,amsthm,mathtools,booktabs}
\usepackage[T1]{fontenc}
\usepackage{lmodern,microtype,needspace}
\usepackage[hidelinks]{hyperref}
\hypersetup{pdftitle={Unique domination in bipartite graphs of order 3 gamma plus 1},pdfauthor={}}
\newtheorem{theorem}{Theorem}[section]
\newtheorem{lemma}[theorem]{Lemma}
\newtheorem{proposition}[theorem]{Proposition}
\newtheorem{corollary}[theorem]{Corollary}
\theoremstyle{definition}
\newtheorem{definition}[theorem]{Definition}
\theoremstyle{remark}
\newtheorem{remark}[theorem]{Remark}
\newcommand{\epn}{\operatorname{epn}}
\newcommand{\e}{e}
\newcommand{\ceil}[1]{\left\lceil#1\right\rceil}
\newcommand{\floor}[1]{\left\lfloor#1\right\rfloor}
\newcommand{\setminusone}[2]{#1\setminus\{#2\}}
\setlength{\emergencystretch}{2em}
\clubpenalty=10000 \widowpenalty=10000
\title{Unique domination in bipartite graphs\\of order $3\gamma+1$}
\author{}
\date{5 October 2026}
\begin{document}
\maketitle
\begin{abstract}
A graph with no isolated vertices and a unique minimum dominating set of
cardinality $\gamma$ has at least $3\gamma$ vertices. We determine the maximum
number of edges in a bipartite graph one vertex above this threshold. For every
$\gamma\ge2$, the maximum at order $3\gamma+1$ is $\gamma(\gamma+7)/2$, and a
connected graph attains it. For $\gamma\ge4$, we characterize the extremal graph
uniquely up to isomorphism. The construction exceeds the corresponding bound
conjectured by Koch and Narayan by $\lfloor\gamma/2\rfloor-1$ edges; in particular,
we recover a previously reported $13$-vertex, $22$-edge counterexample. The proof is elementary. It combines the two
exterior private neighbours forced by uniqueness with replacement arguments
on six and seven vertices, followed by a rigidity analysis of the equality case.
\end{abstract}
\noindent\textbf{Keywords:} unique minimum dominating set; bipartite graph;
extremal size; private neighbour.
\par\smallskip
\noindent\textbf{2020 Mathematics Subject Classification:} 05C69, 05C35.

\section{Introduction}

All graphs in this paper are finite and simple. A set $D\subseteq V(G)$ is
\emph{dominating} if every vertex outside $D$ has a neighbour in $D$.
The domination number $\gamma(G)$ is the minimum cardinality of a dominating
set. A graph has a \emph{unique minimum dominating set} if exactly one dominating
set has cardinality $\gamma(G)$. Throughout, uniqueness refers to minimum
cardinality, not to inclusion-minimality. Write $n(G)=|V(G)|$ and
$\e(G)=|E(G)|$.

Fischermann, Rautenbach and Volkmann~\cite{FRV} investigated the maximum size
of graphs with a unique minimum dominating set. Fraboni and Shank~\cite{FS}
subsequently settled their conjectured size bound when the domination number
is two. Koch and Narayan~\cite{KN} considered the bipartite problem and proposed
a bound depending on the order and the domination number. They proved their
bound for domination number two and at order $n=3\gamma$.
Erlbacher's public artifacts~\cite{E} already give the $13$-vertex, $22$-edge
counterexample, further examples, and finite extremal computations.
The result below recovers that example and proves the sharp boundary maximum
and equality structure for general domination number.
Koch and Narayan also identify the unique extremal graph at
$(n,\gamma)=(10,3)$ in their Figure~1 discussion.

The order restriction is intrinsic. If $D$ is the unique minimum dominating
set of a graph with no isolated vertices, every vertex of $D$ has at least two
exterior private neighbours with respect to $D$. These neighbours are distinct
for distinct members of $D$, and hence $n(G)\ge3\gamma(G)$. We recall the short
argument in Lemma~\ref{lem:private}.

We examine the next possible order. Let $\mathcal B_\gamma$ be the class of
bipartite graphs with no isolated vertices, of order $3\gamma+1$, whose minimum
dominating set is unique and has cardinality $\gamma$. No connectedness
assumption is imposed in the definition.

\Needspace{10\baselineskip}
\begin{theorem}\label{thm:main}
For every integer $\gamma\ge2$,
\[
 \max_{G\in\mathcal B_\gamma}\e(G)=\frac{\gamma(\gamma+7)}2.
\]
The maximum is attained by a connected graph. For $\gamma\ge4$, equality holds
if and only if
\[
 G\cong H\bigl(\ceil{\gamma/2},\floor{\gamma/2}\bigr),
\]
where $H(p,q)$ is the graph defined in Section~\ref{sec:construction}.
\end{theorem}

At order $3\gamma+1$, Conjecture~1 of~\cite{KN} specializes to
\begin{equation}\label{eq:sourcebound}
 \e(G)\le 2\gamma+2ab+2a+1,
 \qquad a=\ceil{\gamma/2},\quad b=\floor{\gamma/2}.
\end{equation}
The summation term in the conjecture is empty at this order. Since
\begin{equation}\label{eq:gap}
 \frac{\gamma(\gamma+7)}2-(2\gamma+2ab+2a+1)=b-1,
\end{equation}
Theorem~\ref{thm:main} has the following consequence.

\begin{corollary}\label{cor:counter}
For every $\gamma\ge4$, there is a connected bipartite graph on $3\gamma+1$
vertices with a unique minimum dominating set of cardinality $\gamma$ that
violates~\eqref{eq:sourcebound}. The excess is $\floor{\gamma/2}-1$.
In particular, at $\gamma=4$ the extremal graph has $13$ vertices and $22$
edges, whereas~\eqref{eq:sourcebound} gives $21$.
\end{corollary}

The additional vertex explains the change from the threshold case $n=3\gamma$.
It may be adjacent to several members of the unique minimum dominating set.
Our extremal graph uses this possibility. The proof of the upper bound keeps
track of whether the additional vertex has one or several such neighbours;
the equality argument then forces the whole construction.

\section{The extremal construction}\label{sec:construction}

\begin{definition}\label{def:H}
For integers $p,q\ge1$, let $H(p,q)$ have pairwise distinct vertices
\[
 X=\{x_i:1\le i\le p\},\qquad Y=\{y_j:1\le j\le q\},
\]
\[
 U_i=\{u_i^0,u_i^1\}\ (1\le i\le p),\qquad
 V_j=\{v_j^0,v_j^1\}\ (1\le j\le q),
\]
and one further vertex $z$. Its edges are exactly
\begin{align*}
 &x_i u_i^\varepsilon &&(1\le i\le p,\ \varepsilon\in\{0,1\}),\\
 &y_j v_j^\varepsilon &&(1\le j\le q,\ \varepsilon\in\{0,1\}),\\
 &u_i^\varepsilon v_j^0 &&(1\le i\le p,\ 1\le j\le q,\ \varepsilon\in\{0,1\}),\\
 &z u_i^\varepsilon &&(1\le i\le p,\ \varepsilon\in\{0,1\}),\\
 &z y_j &&(1\le j\le q).
\end{align*}
\end{definition}

The displayed graph is bipartite, with parts
\begin{equation}\label{eq:bipH}
 X\cup\bigcup_{j=1}^q V_j\cup\{z\}
 \quad\text{and}\quad
 Y\cup\bigcup_{i=1}^p U_i.
\end{equation}
It is connected: $z$ is adjacent to all of $Y\cup\bigcup_iU_i$, each $x_i$
has a neighbour in $U_i$, and each vertex of $V_j$ is adjacent to $y_j$.

\begin{proposition}\label{prop:H}
The graph $H(p,q)$ has order $3(p+q)+1$ and size $2pq+4p+3q$.
Its unique minimum dominating set is $X\cup Y$, and its domination number is
$p+q$.
\end{proposition}

\begin{proof}
The order follows by counting vertices, and the five edge classes in
Definition~\ref{def:H} have respective sizes $2p,2q,2pq,2p,q$.
Clearly $D=X\cup Y$ dominates the graph.

Every dominating set $K$ must meet $N[x_i]=\{x_i,u_i^0,u_i^1\}$ for each $i$.
It must also meet $N[v_j^1]=\{y_j,v_j^1\}$ for each $j$, since $v_j^1$ is a
leaf. These $p+q$ sets are pairwise disjoint, so $|K|\ge p+q$.
This proves $\gamma(H(p,q))=p+q$.

Suppose now that $|K|=p+q$. It chooses exactly one vertex from each of these
sets and no other vertex. In particular, $z$ and all $v_j^0$ are absent.
If $K$ contains $u_i^0$ rather than $x_i$, then $u_i^1$ is undominated:
its neighbours are $x_i,z$ and the vertices $v_j^0$. The same argument with
the superscripts reversed excludes $u_i^1$. Thus every $x_i$ belongs to $K$
and no vertex of any $U_i$ belongs to $K$.
Finally, choosing $v_j^1$ rather than $y_j$ leaves $v_j^0$ undominated,
because its neighbours are $y_j$ and the vertices in the sets $U_i$.
Hence $K=X\cup Y$.
\end{proof}

For $a=\ceil{\gamma/2}$ and $b=\floor{\gamma/2}$, a direct calculation gives
\begin{equation}\label{eq:attain}
 \e(H(a,b))=2ab+4a+3b=\frac{\gamma(\gamma+7)}2.
\end{equation}
Proposition~\ref{prop:H} therefore supplies the lower bound in
Theorem~\ref{thm:main}, including connectedness.

\section{Private neighbours and local replacements}

For a dominating set $D$ and $d\in D$, define
\[
 \epn(d,D)=\{u\in V(G)\setminus D:N(u)\cap D=\{d\}\}.
\]
We use closed neighbourhoods only when square brackets are displayed.

\begin{lemma}\label{lem:private}
If $G$ has no isolated vertices and $D$ is its unique minimum dominating set,
then $|\epn(d,D)|\ge2$ for every $d\in D$.
\end{lemma}

\begin{proof}
If $\epn(d,D)=\{u\}$, the set $(D\setminus\{d\})\cup\{u\}$ dominates $G$:
$u$ dominates itself and $d$, while every other vertex outside $D$ retains a
neighbour in $D\setminus\{d\}$. This contradicts uniqueness.

If $\epn(d,D)=\varnothing$ and $d$ has a neighbour in $D$, then
$D\setminus\{d\}$ is a smaller dominating set. Otherwise $d$ has no neighbour
in $D$. The no-isolate assumption gives a neighbour $u$ outside $D$, and
$(D\setminus\{d\})\cup\{u\}$ is again a different dominating set of the same
cardinality. All cases are impossible.
\end{proof}

The next two lemmas concern induced subgraphs that arise after two members
of $D$ are removed. They include every possible edge between the two centers.
The six-vertex argument is the local estimate used in the proof of
Theorem~12 of~\cite{KN}; we record its equality patterns as well.

\begin{lemma}\label{lem:six}
Let $F$ be a bipartite graph with parts $\{x\}\cup V$ and $\{y\}\cup U$,
where $|U|=|V|=2$. Assume that $x$ is adjacent to both vertices of $U$ and
$y$ to both vertices of $V$. If $\{x,y\}$ is the unique dominating set of
cardinality at most two, then at most two edges occur among $xy$ and the
$U$--$V$ edges. If exactly two occur, either $xy$ is absent and the
$U$--$V$ edges form a two-edge star, or $xy$ is present and there is exactly
one $U$--$V$ edge.
\end{lemma}

\begin{proof}
Two disjoint $U$--$V$ edges give an alternative dominating pair: select
opposite endpoints, one from each edge. If $xy$ is absent, three
$U$--$V$ edges contain two disjoint edges. Thus there are at most two, and a
two-edge configuration must be a star.

Suppose $xy$ is present. Two $U$--$V$ edges either are disjoint or form a
star. A star centered at $u\in U$ gives the dominating pair $\{x,u\}$;
a star centered at $v\in V$ gives $\{y,v\}$. Consequently at most one
$U$--$V$ edge can accompany $xy$.
\end{proof}

\begin{lemma}\label{lem:seven}
Let $F$ have the same form as in Lemma~\ref{lem:six}, except that $|U|=2$
and $|V|=3$. If $\{x,y\}$ is the unique dominating set of cardinality at
most two, then at most four edges occur among $xy$ and the $U$--$V$ edges.
\end{lemma}

\begin{proof}
If $xy$ is absent and at least five $U$--$V$ edges occur, one vertex
$u\in U$ is adjacent to all three vertices of $V$, and the other vertex
$u'$ is adjacent to at least two. Choose $v\in V$ adjacent to $u'$.
The pair $\{u,v\}$ dominates $F$, a contradiction.

If $xy$ is present, any vertex $v\in V$ adjacent to both vertices of $U$
gives the alternative dominating pair $\{y,v\}$. Every vertex of $V$
therefore has at most one neighbour in $U$. There are at most three
$U$--$V$ edges, and at most four edges when $xy$ is included.
\end{proof}

These are direct combinatorial proofs. The exhaustive checks summarized in
Appendix~\ref{app:checks} are supplementary and are not used as hypotheses in
what follows.

\section{The sharp upper bound}\label{sec:upper}

Let $G\in\mathcal B_\gamma$, and let $D$ be its unique minimum dominating set.
By Lemma~\ref{lem:private}, we may select two exterior private neighbours for
each vertex of $D$. The selected pairs are disjoint and account for $2\gamma$
vertices. Since $n(G)=3\gamma+1$, exactly one further vertex $z$ remains.

Choose the names of the bipartition classes so that $z\in A$, and write
\[
 D\cap A=\{x_1,\ldots,x_p\},\qquad
 D\cap B=\{y_1,\ldots,y_q\},\qquad p+q=\gamma.
\]
Let $U_i\subseteq B$ be the selected private pair of $x_i$, and
$V_j\subseteq A$ the selected private pair of $y_j$. Set
\[
 J=N(z)\cap D\subseteq\{y_1,\ldots,y_q\},\qquad t=|J|.
\]
Because $D$ dominates $z$, we have $t\ge1$, and in particular $q\ge1$.

Privacy and bipartiteness give a complete list of possible edges:
\begin{enumerate}
\item the $2\gamma$ selected center-to-private edges;
\item $x_iy_j$ and the edges between $U_i$ and $V_j$, for each pair $(i,j)$;
\item the edges from $z$ to $J$ and to $\bigcup_iU_i$.
\end{enumerate}
In particular, edges inside $D$ have not been discarded: they are precisely
the possible edges $x_iy_j$ in the second class.

\subsection{The additional vertex has several dominators}

Suppose $t\ge2$. For each $i,j$, consider the induced six-vertex graph
\[
 G[K_{ij}],\qquad K_{ij}=\{x_i,y_j\}\cup U_i\cup V_j.
\]
The set $D\setminus\{x_i,y_j\}$ dominates every vertex outside $K_{ij}$.
Each selected private vertex outside the cell retains its own center;
the vertex $z$ retains a neighbour in $J$ because only one $y$-center has
been removed.

If $G[K_{ij}]$ had a dominating set $S$ of cardinality at most two other
than $\{x_i,y_j\}$, then
\[
 (D\setminus\{x_i,y_j\})\cup S
\]
would be a dominating set of $G$ of cardinality at most $\gamma$, distinct
from $D$ when its cardinality is $\gamma$. This contradicts minimum
cardinality or uniqueness. Thus Lemma~\ref{lem:six} applies to every cell.
The optional edge sets of these cells are disjoint, and
$\deg(z)\le2p+q$. Consequently
\begin{equation}\label{eq:multiple}
 \e(G)\le 2\gamma+2pq+2p+q.
\end{equation}
When $p=0$, the same bound follows directly from the edge list, with an
empty sum over cells.

\subsection{The additional vertex has one dominator}

Suppose $J=\{y_{j_0}\}$. For every $j\ne j_0$, the preceding six-vertex
replacement still applies, since $y_{j_0}$ is retained and dominates $z$.
For each $i$ consider instead
\[
 L_i=\{x_i,y_{j_0}\}\cup U_i\cup V_{j_0}\cup\{z\}.
\]
In $G[L_i]$, the center $y_{j_0}$ has the three private vertices
$V_{j_0}\cup\{z\}$. Every vertex outside $L_i$ is dominated by
$D\setminus\{x_i,y_{j_0}\}$, so the same replacement argument proves that
Lemma~\ref{lem:seven} applies.

The optional edges of $G[L_i]$ are $x_i y_{j_0}$ and the edges from $U_i$
to $V_{j_0}\cup\{z\}$. These edge sets are disjoint for different $i$.
The edge $zy_{j_0}$ is counted once separately. Therefore
\begin{align}
 \e(G)&\le 2\gamma+1+2p(q-1)+4p\notag\\
       &=2\gamma+1+2pq+2p\notag\\
       &\le2\gamma+2pq+2p+q.\label{eq:single}
\end{align}
If $p=0$, only the $2\gamma$ fixed private edges and $zy_{j_0}$ occur,
so the displayed inequalities remain valid.

\subsection{Optimizing the bipartition}

Put $d=p-q$. Since $p+q=\gamma$, both cases yield
\begin{equation}\label{eq:opt}
 \e(G)\le 2\gamma+2pq+2p+q
 =\frac{\gamma(\gamma+7)}2-\frac{d(d-1)}2
 \le\frac{\gamma(\gamma+7)}2.
\end{equation}
The final inequality holds because $d$ is an integer. Equality in the
numerical optimization requires $d\in\{0,1\}$. The parity of $\gamma$
then forces
\begin{equation}\label{eq:balanced}
 p=\ceil{\gamma/2},\qquad q=\floor{\gamma/2}.
\end{equation}
Together with~\eqref{eq:attain}, this proves the extremal-value assertion of
Theorem~\ref{thm:main} for every $\gamma\ge2$.

\section{Rigidity at equality}\label{sec:rigidity}

Suppose now that $\gamma\ge4$ and
$\e(G)=\gamma(\gamma+7)/2$. Equation~\eqref{eq:balanced} gives
$q=\floor{\gamma/2}\ge2$. The final inequality in~\eqref{eq:single}
loses $q-1$ edges, so the case $t=1$ cannot give equality.
Every inequality in~\eqref{eq:multiple} must therefore be an equality.
It follows that $z$ is adjacent to every $y_j$ and to every vertex of
$\bigcup_iU_i$, and each six-vertex cell has exactly two optional edges.

We first show that $D$ is independent. If $x_iy_j$ were an edge, then
\[
 (D\setminus\{x_i\})\cup\{z\}
\]
would be a different dominating set of cardinality $\gamma$.
Indeed, $x_i$ is dominated by the retained $y_j$, the pair $U_i$ is
dominated by $z$, and every other private pair retains its center.
Thus no edge $x_iy_j$ occurs.

Lemma~\ref{lem:six} now shows that the two edges between $U_i$ and $V_j$
form a star. This star cannot be centered at $u\in U_i$, because
\[
 (D\setminus\{x_i,y_j\})\cup\{u,z\}
\]
would dominate $G$. The vertex $u$ dominates $x_i$ and both members of
$V_j$, while $z$ dominates $y_j$ and both members of $U_i$.
All other private pairs retain their centers. Therefore both members of
$U_i$ are adjacent to exactly one vertex of $V_j$; denote it by $v_j^{(i)}$.

Fix $j$. Suppose that the selected vertex $v_j^{(i)}$ is not constant as
$i$ varies. Choose one vertex $u_i\in U_i$ for every $i$, and define
\[
 K=\{z\}\cup\{u_i:1\le i\le p\}
       \cup\bigl((D\cap B)\setminus\{y_j\}\bigr).
\]
Then $|K|=1+p+(q-1)=\gamma$. Every $x_i$ is dominated by $u_i$;
every vertex in $\bigcup_iU_i$ and every $y$-center is dominated by $z$.
The two vertices of $V_j$ are dominated because both possible choices of
$v_j^{(i)}$ occur. Every other pair $V_\ell$ retains $y_\ell$.
Hence $K$ is a second minimum dominating set, a contradiction.

For each $j$, rename the members of $V_j$ so that their common selected
vertex is $v_j^0$. The complete edge list is now exactly that of $H(p,q)$.
Equation~\eqref{eq:balanced} proves the equality assertion in
Theorem~\ref{thm:main}.

\section{Consequences and further questions}

The extremal values at $\gamma=2,3$ are $9$ and $15$; the graphs
$H(1,1)$ and $H(2,1)$ attain them. The rigidity assertion proved above is
restricted to $\gamma\ge4$. Neither the theorem nor Corollary~\ref{cor:counter}
asserts that the $13$-vertex example is the smallest counterexample over all
orders in the conjecture of~\cite{KN}.

For $\gamma\ge4$, the additional vertex of the extremal graph has
$q=\floor{\gamma/2}\ge2$ neighbours in the unique minimum dominating set.
Consequently this set is not an efficient dominating set. The shared vertex
is forced at equality and is responsible for the additional edges over
\eqref{eq:sourcebound}. At $n=3\gamma$, there is no vertex available for this
role.

A natural next problem is to determine the extremal size at $n=3\gamma+r$
for a fixed $r\ge2$, together with the equality cases. Selecting two private
neighbours for each dominator then leaves $r$ vertices. Replacing two centers
may leave some of those vertices undominated, so their incidences with the
dominating set must be retained in any local reduction. The exact threshold
at which new patterns arise is not determined here.

\appendix
\section{Supplementary exact checks}\label{app:checks}

The proofs in the main text do not depend on a computer enumeration.
Two separate implementations were nevertheless used to check the local
lemmas, the explicit counterexample, and the first values of the construction.
All tests use finite sets or integer bit masks.

For the six-vertex cell, the fixed edges are $xu_0,xu_1,yv_0,yv_1$.
The optional edges, in bit order, are
\[
 xy,\ u_0v_0,\ u_0v_1,\ u_1v_0,\ u_1v_1.
\]
For the seven-vertex cell, the fixed edges are $xu_0,xu_1,yv_0,yv_1,yv_2$,
and the optional edges are
\[
 xy,\ u_0v_0,\ u_0v_1,\ u_0v_2,\ u_1v_0,\ u_1v_1,\ u_1v_2.
\]
Each implementation considers every optional-edge subset and tests every
vertex subset of cardinality at most two for domination. Table~\ref{tab:cells}
records the results. ``Admissible'' means that $\{x,y\}$ is the only dominating
set of cardinality at most two. Decimal masks use the first listed edge as
the least significant bit.

\begin{table}[htbp]
\centering
\small
\begin{tabular}{lrrrl}
\toprule
Cell & Patterns & Admissible & Maximum & Maximizing masks\\
\midrule
Six vertices & 32 & 14 & 2 & 3, 5, 6, 9, 10, 17, 20, 24\\
Seven vertices & 128 & 58 & 4 & 54, 90, 108\\
\bottomrule
\end{tabular}
\caption{Exhaustive local checks. The maximum counts optional edges.}
\label{tab:cells}
\end{table}

A separate verifier reads the explicit $13$-vertex edge list and checks all
$2^{13}=8192$ vertex subsets. It finds no dominating set of size at most three
and exactly one of size four, namely $X\cup Y$. It also checks connectedness,
bipartiteness, simplicity and the edge count $22$. The counts of dominating
sets, by cardinality $0,1,\ldots,13$, are
\[
 (0,0,0,0,1,127,492,855,876,586,262,76,13,1).
\]
The first balanced construction values are recorded in
Table~\ref{tab:small}; all subsets of cardinality at most $\gamma$ were checked.

\begin{table}[htbp]
\centering
\begin{tabular}{rrrrr}
\toprule
$\gamma$ & $(p,q)$ & $n$ & $\e(H(p,q))$ & Minimum dominating sets\\
\midrule
2 & $(1,1)$ & 7 & 9 & 1\\
3 & $(2,1)$ & 10 & 15 & 1\\
4 & $(2,2)$ & 13 & 22 & 1\\
5 & $(3,2)$ & 16 & 30 & 1\\
6 & $(3,3)$ & 19 & 39 & 1\\
\bottomrule
\end{tabular}
\caption{Small instances of the connected extremal family.}
\label{tab:small}
\end{table}

The accompanying files \path{check_local_cells.py} and
\path{local_cells_certificate.json} give all local patterns and their
dominating pairs. The explicit example is in \path{candidate13.json};
\path{check_counterexample_independent.py} verifies it from that edge list.
A second implementation, \path{audit/independent_check.py}, independently
repeats the local checks and tests all $4096$ saturated assemblies at $p=q=2$.
Exactly four labelled assemblies retain a unique minimum dominating set;
they correspond to the independent choice of which member of each $V_j$
is named $v_j^0$, and are isomorphic to $H(2,2)$.

\begin{thebibliography}{9}
\bibitem{FRV}
M.~Fischermann, D.~Rautenbach and L.~Volkmann,
\emph{Maximum graphs with a unique minimum dominating set},
Discrete Mathematics \textbf{260} (2003), 197--203.

\bibitem{FS}
M.~Fraboni and N.~Shank,
\emph{Maximum graphs with unique minimum dominating set of size two},
Australasian Journal of Combinatorics \textbf{46} (2010), 91--99.
\url{https://ajc.maths.uq.edu.au/pdf/46/ajc_v46_p091.pdf}.

\bibitem{E}
J.~Erlbacher,
\emph{Demonstrandum---wave-1 verified artifacts: Koch--Narayan counterexample package},
public research artifacts, 2026.
\href{https://github.com/demonstrandum-research/artifacts/blob/94db9ed50d48a57aae5ccb72e6a95a2b8f8f39d3/problems/p2-factory/kills/koch-narayan/WRITEUP.md}{Pinned write-up, commit \texttt{94db9ed50d48} (13 June 2026)}.

\bibitem{KN}
G.~Koch and D.~Narayan,
\emph{Maximal bipartite graphs with a unique minimum dominating set},
arXiv:2511.01719v1, 3 November 2025.
\url{https://arxiv.org/abs/2511.01719v1}.
\end{thebibliography}
\end{document}
